\section{Evaluation Setup}
\label{sec:setup}

This section fixes the protocol: how the benchmark is split, what is
measured, which systems are compared, and how significance is decided.

\subsection{Time-based split}
\label{sec:setup:protocol}

Each project's instances are ordered by the time of their fix commit. The
earliest 30\% form the development split, and the later 70\% form the
evaluation split. All tuning, which in \sys{} means only the fusion
weights of \cref{sec:approach:fusion}, happens on the development split.
Every number reported in \cref{sec:results} comes from the evaluation
split. We chose a time split rather than a random one because it mirrors
deployment: a system is configured on the project's past and then meets
the project's future. A random split would let instances from the future
inform the configuration used on the past.

The time split has a side effect worth recording, because any system that
uses historical memory will meet it. The development split holds the
oldest issues, and the oldest issues have the least history behind them.
Tuned on that slice, the fusion weights undervalue memory relative to its
worth on later data. We report this as a property of the protocol and
keep the split anyway; the alternative, tuning on recent data and
evaluating on older data, would be the deployment scenario backwards.

\subsection{Metrics}
\label{sec:setup:metrics}

We report the challenge's metrics at the file level: Acc@K, the fraction
of instances with at least one gold file in the top $K$; Recall@K, the
average fraction of an instance's gold files found in the top $K$; and
MRR, the mean reciprocal rank of the first gold file. We use
$K \in \{1, 5, 10\}$ and treat Acc@5 as primary, since five files is a
realistic number for an engineer to open. At the module level we report
Acc@K and MRR for the module ranking defined in
\cref{sec:approach:fusion}; the function-level stage of
\cref{sec:results:functions} reuses the same metrics over function
candidates.

\subsection{Compared systems}
\label{sec:setup:systems}

Each signal from \cref{sec:approach} runs standalone as a baseline:
trace mentions (S1), path retrieval (S2), content retrieval (S3), and
memory alone (S4). S1 represents what an engineer gets by grepping the
report. On top of these we evaluate the fused system in two forms,
unweighted fusion and regime-conditioned weighted fusion, and each form
both with and without the memory signal. The without-memory variants are
tuned with the same grid and the same development split as the full
system, so the memory ablation compares two equally tuned systems rather
than a tuned system against a crippled one.

We additionally re-implement the classic IRBL baselines, since the
original tools are Java-only. \textbf{rVSM} scores every file of the full
snapshot with a length-favoring tf--idf cosine. \textbf{BugLocator}
combines rVSM with SimiScore, the 2012 similar-report
signal~\cite{zhou2012buglocator}: every past fixed report votes for its
fix's files with weight $\mathrm{sim}(q,b)/n_b$; the blend weight uses
the replication-standard $\alpha = 0.3$, with
$\alpha \in \{0.2, 0.3, 0.5\}$ tuned on the development split as a
variant. \textbf{BRTracer-style} adds the stack-trace boost of
BRTracer~\cite{wong2014brtracer} ($+1/r$ for the file at trace rank
$r \le 10$); we omit the original's file segmentation and state so. Unless noted otherwise, their SimiScore runs
under the SAME firewall as our memory signal, so the comparison is honest
on both sides; \cref{sec:results:leakage} measures what happens without
it.

\subsection{Statistical treatment}
\label{sec:setup:stats}

Every comparison in \cref{sec:results} is tested with a paired bootstrap
over instances: 2{,}000 resamples, identical resample indices for all
comparisons within a population, two-sided $p$-values from the bootstrap
distribution, and 95\% confidence intervals. The paper reports two test
families, each corrected with the Benjamini--Hochberg procedure over the
whole family: 64 tests for the main comparisons (four comparisons, two
metrics, two populations, four projects) and 96 tests for the baseline
and leakage experiment of \cref{sec:results:leakage} (six comparisons,
same grid). We call a result significant only when the corrected
$q < 0.05$. The pipeline itself is deterministic from end to end: fixed
seeds, greedy decisions, and no GPU training. Re-running every experiment
reproduces every table in this paper bit for bit, and the released
artifact contains a script that traces each reported number to the result
file it came from.
